%%%%%%%%%%  scan idx 572  |  folio 559  %%%%%%%%%%
\SGAPrintPage{559}%

power. For example, on the topos \((\Ens)\) ringed by the field \(\mathbb Q\), the representation \(G_m\to G_m\) defined by
\[
 r\rightsquigarrow|r|\qquad(r\in\mathbb Q^*)
\]
is not a power. In particular, 6.2.3 does not remain valid for an arbitrary locally ringed topos \(X\) and an arbitrary representation \(\alpha:G\to G_m\).

\sgaStableTarget{sga6.e10.para.6-5-3}\textbf{6.5.3.} On the other hand, 6.2.3 remains valid for the representations associated with formulas (6.2.8), (6.2.9) and (6.2.10) on an arbitrary ringed topos \(X=({C},\underline{O})\), because these representations are ``algebraic'' in the following sense:

Let \(\underline{e}\) be the final object of \(X\), put \(S=\Spec(\underline{O}(\underline{e}))\), and consider the functor
\[
 u:{C}\longrightarrow(\Sch/S)
\]
which assigns to each \(Y\in\operatorname{Ob}({C})\) the affine scheme \(\Spec(\underline{O}(Y))\) over \(S\). Then, for every positive integer \(n\), the presheaf \(\operatorname{G\ell}(n)_X\) on \({C}\) is the composite with \(u\) of the corresponding presheaf \(\operatorname{G\ell}(n)_S\) on \((\Sch/S)\). A representation on \(X\)
\sgaStableTarget{sga6.e10.eq.6-5-4}\begin{equation}\tag{6.5.4}
 \alpha:\operatorname{G\ell}(n)\longrightarrow G_m
\end{equation}
is called \emph{algebraic} if it arises via \(u\) from a representation
\sgaStableTarget{sga6.e10.eq.6-5-5}\begin{equation}\tag{6.5.5}
 \alpha_S:\operatorname{G\ell}(n)_S\longrightarrow G_{m,S}
\end{equation}
of presheaves of groups on \((\Sch/S)\), i.e., if one has
\[
 \alpha(s)=\alpha_S(s)
\]
for every \(Y\in\operatorname{Ob}({C})\), \(s\in\operatorname{G\ell}(n)(Y)=\operatorname{G\ell}(n)_S(u(Y))\).

The representation \(\alpha\) in (6.5.4) is, of course, completely determined by \(\alpha_S\). It follows immediately that (6.2.3) remains valid in an arbitrary ringed topos when \(\alpha\) is algebraic. In particular, formulas (6.2.8), (6.2.9) and (6.2.10) therefore remain valid on an arbitrary ringed topos.

%%%%%%%%%%  scan idx 573  |  folio 560  %%%%%%%%%%
\SGAPrintPage{560}%

\sgaStableTarget{sga6.e10.sec.7}\section*{7. Appendix: Specialization in Intersection Theory}
\begin{center}
\underline{by A. Grothendieck}
\end{center}

\sgaStableTarget{sga6.e10.lem.7-1}\textbf{Lemma 7.1.} \uline{Let} \(i:Y\to X\) \uline{be a regular closed immersion} (VII 1.4)\uline{, of bounded codimension, so that} \(i\) \uline{is a perfect morphism} (VII 1.9)\uline{, and one can therefore define the morphism} (IV 2.12)
\[
 i^*:K_\bullet(X)\longrightarrow K_\bullet(Y).
\]
\uline{Let, on the other hand,} \(\mathbb N\) \uline{be the conormal sheaf of} \(i\)\uline{, which is a locally free sheaf of finite type on} \(Y\) (VII 1.6)\uline{, so that one can form}
\[
 \lambda_{-1}(\mathbb N)\in K^{\bullet}(Y).
\]
\uline{Consider also the direct-image homomorphism} (IV 2.11)
\[
 i_*:K_\bullet(Y)\longrightarrow K_\bullet(X).
\]
\uline{Then one has the formula}
\[
 i^*i_*(x)=x\lambda_{-1}(\mathbb N)\qquad\text{for } x\in K_\bullet(Y),
\]
\uline{where the product on the right is taken with respect to the} \(K^{\bullet}(Y)\)\uline{-module structure on} \(K_\bullet(Y)\) \uline{defined in} (IV 2.10).

The proof is, essentially without change, the same as that of the same formula, written for \(K^{\bullet}(X)\), \(K^{\bullet}(Y)\), established in (VII 2.7).

\sgaStableTarget{sga6.e10.cor.7-2}\textbf{Corollary 7.2.} \uline{Let} \(f:X\to S\) \uline{be a flat morphism, with} \(S\) \uline{a trait, i.e., the spectrum of a discrete valuation ring} \(A\), \(s\) \uline{the closed point of} \(S\), \(t\) \uline{its generic point, and} \(i:X_s\to X\) \uline{the canonical inclusion. Then} \(i\) \uline{has finite Tor-dimension, so} (IV 2.12) \uline{one can define}
\[
 i^*:K_\bullet(X)\longrightarrow K_\bullet(X_s),
\]

%%%%%%%%%%  scan idx 574  |  folio 561  %%%%%%%%%%
\SGAPrintPage{561}%

\uline{on the other hand}, (IV 2.11) \uline{one has a homomorphism}
\[
 i_*:K_\bullet(X_s)\longrightarrow K_\bullet(X).
\]
\uline{Then}, \uline{the composite}
\[
 i^*i_*:K_\bullet(X_s)\longrightarrow K_\bullet(X)\longrightarrow K_\bullet(X_s)
\]
\uline{is zero}.

Indeed, \(i:X_s\to X\) is a regular immersion of codimension \(1\) whose conormal sheaf \(\mathbb N\) is \emph{free} of rank \(1\), since \(X_s\) can be described by a global equation \(f^*(g)=0\) (where \(g\) is a uniformizer of \(A\)). Thus one has \(\lambda_{-1}(\mathbb N)=0\), and the conclusion follows from 7.1.

Throughout what follows, we shall suppose that \(X\) is \emph{noetherian} and \emph{flat over} \(S\). Let us now apply the exact homotopy sequence (IX 1.1) to the following situation: \(X\), with the open subscheme \(X_t\) and complementary closed subscheme \(X_s\):
\[
K_\bullet(X_s)\xrightarrow{i_*}K_\bullet(X)\xrightarrow{j^*}K_\bullet(X_t)\longrightarrow0,
\]
where \(j:X_t\to X\) is the inclusion, which is an open immersion. Using 7.2, one obtains:

\sgaStableTarget{sga6.e10.prop.7-3}\textbf{Proposition 7.3.} \uline{With the notation of} 7.2, \(X\) \uline{being assumed noetherian}, \uline{there is a unique group homomorphism}
\sgaStableTarget{sga6.e10.eq.7-3-1}\begin{equation}\tag{7.3.1}
 \sigma_\bullet:K_\bullet(X_t)\longrightarrow K_\bullet(X_s)
\end{equation}
(\uline{called the} ``specialization homomorphism'') \uline{making the diagram commutative}
\sgaStableTarget{sga6.e10.eq.7-3-2}\begin{equation}\tag{7.3.2}
\begin{tikzcd}[column sep=large,row sep=large]
K_\bullet(X) \arrow[r,"j^*"] \arrow[d,"i^*"'] & K_\bullet(X_t) \arrow[dl,"\sigma_\bullet"]\\
K_\bullet(X_s) & {.}
\end{tikzcd}
\end{equation}


% Local macro additions for Exposé X continuation
\providecommand{\Kalg}{K_{\mathrm{alg}}}
\providecommand{\Gralg}{\operatorname{Gr}_{\mathrm{alg}}}
\providecommand{\Grnum}{\operatorname{Gr}_{\mathrm{num}}}
\providecommand{\Knum}{K_{\mathrm{num}}}
\providecommand{\Cl}{\operatorname{Cl}}
\providecommand{\Ch}{\operatorname{Ch}}
\providecommand{\Nm}{\operatorname{N}}
\providecommand{\colim}{\varinjlim}
\providecommand{\Hens}{\operatorname{Hens}}
\providecommand{\Spec}{\operatorname{Spec}}
\providecommand{\coker}{\operatorname{coker}}


%%%%%%%%%%  scan idx 575  |  folio 562  %%%%%%%%%%
\SGAPrintPage{562}%

\noindent \uline{When the two sides of} (7.3.1) \uline{are regarded as modules over} \(K^{\bullet}(X)\) \uline{by means of the ring homomorphisms} \(K^{\bullet}(X)\xrightarrow{j^*}K^{\bullet}(X_t)\) \uline{and} \(K^{\bullet}(X)\xrightarrow{i^*}K^{\bullet}(X_s)\) (IV 2.7 b) \uline{and the module structures on} \(K_\bullet(X_t)\) \uline{and} \(K_\bullet(X_s)\) \uline{over} \(K^{\bullet}(X_t)\) \uline{and} \(K^{\bullet}(X_s)\) \uline{respectively} (IV 2.10), \uline{the homomorphism} \(\sigma_\bullet\) \uline{is a homomorphism of} \(K^{\bullet}(X)\)-\uline{modules}.

The last assertion follows from the fact that \(i^*:K_\bullet(X)\to K_\bullet(X_s)\) is a homomorphism of \(K^{\bullet}(X)\)-modules, as is the epimorphism \(j^*:K_\bullet(X)\longrightarrow K_\bullet(X_t)\), by (IV 2.12.2).

We shall review other properties of the specialization morphism.

\sgaStableTarget{sga6.e10.sec.7.4}\subsection*{7.4. Compatibility with filtrations.}
Endow the \(K_\bullet\) with the filtration \(\operatorname{Fil}_\bullet K_\bullet\) by the dimension of supports (1.1.1). Then, when \(f\) is of finite type, more generally when \(X\) is catenary, \emph{the specialization homomorphism is compatible with these filtrations}, and consequently induces a homomorphism (also called the specialization homomorphism):
\sgaStableTarget{sga6.e10.eq.7-4-1}\begin{equation}\tag{7.4.1}
  \sigma_\bullet:\Gr_\bullet(X_t)\longrightarrow \Gr_\bullet(X_s).
\end{equation}

To see this, by 1.1.4 one is reduced to proving that, when \(Y_t\) is an integral subscheme of dimension \(d\) of \(X_t\), \(\sigma_\bullet(\cl_\bullet(Y_t))\) has filtration \(\leq d\). Now let \(Y\) be the schematic closure of \(Y_t\) in \(X\), so that \(Y\) is flat over \(S\) (EGA IV 2.8.1). It follows readily that \(i^*(\cl_\bullet(Y))=\cl_\bullet(Y_s)\), where the two \(\cl_\bullet\) are taken in \(K_\bullet(X)\) and \(K_\bullet(X_s)\) respectively, whence \(\sigma_\bullet(\cl_\bullet(Y_t))=\cl_\bullet(Y_s)\). One checks readily (using the fact that \(X\) is catenary, hence \(Y\) is catenary) that \(\dim Y_s\leq \dim Y_t\); when \(f\) is of finite type, this follows at once, by reduction to the case of integral \(X\), from EGA IV 7.1.13. This gives the asserted conclusion.

%%%%%%%%%%  scan idx 576  |  folio 563  %%%%%%%%%%
\SGAPrintPage{563}%

\sgaStableTarget{sga6.e10.sec.7.5}\subsection*{7.5. ``Conservation of number''.}
Suppose now that \(f\) is proper, so that \(X_t\) and \(X_s\) are proper over \(k(t)\) and \(k(s)\) respectively, and one has the homomorphisms of 2.3:
\[
  \chi_t:K_\bullet(X_t)\longrightarrow \mathbb Z,\qquad
  \chi_s:K_\bullet(X_s)\longrightarrow \mathbb Z .
\]
Then one has the relation
\sgaStableTarget{sga6.e10.eq.7-5-1}\begin{equation}\tag{7.5.1}
  \chi_s(\sigma_\bullet(x))=\chi_t(x)\qquad\text{for } x\in K_\bullet(X_t)
\end{equation}
(the ``conservation-of-number theorem''); in other words, the diagram
\sgaStableTarget{sga6.e10.eq.7-5-2}\begin{equation}\tag{7.5.2}
\begin{tikzcd}[column sep=large,row sep=large]
K_\bullet(X_t) \arrow[r,"\sigma_\bullet"] \arrow[dr,"\chi_t"'] &
K_\bullet(X_s) \arrow[d,"\chi_s"]\\
& \mathbb Z .
\end{tikzcd}
\end{equation}
commutes. To prove (7.5.1), one is again reduced to the case where \(x=\cl_\bullet(Y_t)\), with \(Y_t\) an integral closed subscheme of \(X_t\), and, with the notation of the proof of 7.4, one is reduced to proving that
\[
  \chi(Y_s)=\chi(Y_t),
\]
which is well known (EGA III 7.9.4, or Exposé III 5.7).

\sgaStableTarget{sga6.e10.sec.7.6}\subsection*{7.6.}
It is immediate that
\sgaStableTarget{sga6.e10.eq.7-6-1}\begin{equation}\tag{7.6.1}
  \sigma_\bullet\bigl(\cl_\bullet(\underline{O}_{X_t})\bigr)=\cl_\bullet(\underline{O}_{X_s}),
\end{equation}
which, taking account of the \(K^{\bullet}(X)\)-linearity of \(\sigma_\bullet\), implies that the following diagram is commutative:
\sgaStableTarget{sga6.e10.eq.7-6-2}\begin{equation}\tag{7.6.2}
\begin{tikzcd}[column sep=large,row sep=large]
& K^{\bullet}(X) \arrow[dl,"j^*"'] \arrow[dr,"i^*"] &\\
K^{\bullet}(X_t) \arrow[d] && K^{\bullet}(X_s) \arrow[d]\\
K_\bullet(X_t) \arrow[rr,"\sigma_\bullet"] && K_\bullet(X_s),
\end{tikzcd}
\end{equation}
